\section{课程导读}
%! src=1 kind=orig title=课件封面与整理说明
\begin{frame}[t,shrink]{\src{1}{课件封面与整理说明}}
\begin{columns}[T,onlytextwidth]
\begin{column}{0.57\textwidth}
  {\small\bfseries 这份幻灯片的三个层次}\par\vspace{1.2mm}
  \begin{itemize}
    \item \srctag\ \textbf{原页总结}：逐页对应原课件，保留定义、流程、公式与数据；标题中的“原第 \(XX\) 页”即原 PDF 页码。
    \item \addtag\ \textbf{补充推导}：补充物理含义、变量单位、公式条件、必要推导与易混点。
    \item \readtag\ \textbf{读图提示}：说明流程图、焓熵图、速度三角形和结构图应先看什么、再看什么。
  \end{itemize}
  \vspace{1.0mm}
  \begin{keybox}{建议阅读顺序}
  \small 原页总结 $\rightarrow$ 补充推导 $\rightarrow$ 合上页面复述“设备作用—介质流向—能量边界—公式条件”。
  \end{keybox}
\end{column}
\begin{column}{0.40\textwidth}
  \begin{plainbox}
  {\small\bfseries 原课件信息}\par\vspace{1.2mm}
  \begin{tabularx}{\linewidth}{@{}>{\bfseries}p{15mm}X@{}}
    课程 & 热能动力工程\\[0.8mm]
    教师 & 于立军、陈金伟\\[0.8mm]
    篇幅 & 原 PDF 共 72 页\\[0.8mm]
    章节 & 能源用能、锅炉、汽轮机、燃气轮机、内燃机\\[0.8mm]
    底稿 & Marp 逐页总结\\[0.8mm]
    索引 & 原页—新页页码对照表
  \end{tabularx}
  \end{plainbox}
  \vspace{1.0mm}
  {\fontsize{7.2}{9}\selectfont\textcolor{black!58}{说明：课件页码与 PDF 物理页序一致；本整理不修改原 PDF。}}
\end{column}
\end{columns}
\vspace{1mm}
\begin{plainbox}
\small 全稿共 103 页：封面与课程地图 2 页 + 原页总结 72 页 + 补充 29 页。页码索引见 \texttt{slide-index.tsv}。
\end{plainbox}
\end{frame}

%! src=0 kind=nav title=课程地图
\begin{frame}[t,shrink]{课程地图｜五个单元与一条能量链}
\begin{center}
\begin{tikzpicture}[x=1mm,y=1mm]
  \draw[cprimary!25,line width=1.5pt] (8,0) -- (152,0);
  \foreach \x/\n/\t/\p in {%
    18/01/能源与科学用能/{原页 2--7},
    50/02/锅炉/{原页 8--24},
    82/03/汽轮机/{原页 25--46},
    114/04/燃气轮机/{原页 47--59},
    144/05/内燃机/{原页 60--72}}{
    \fill[cprimary] (\x,0) circle (2.2);
    \node[align=center,font=\fontsize{7.0}{9}\selectfont\bfseries,text=black!80] at (\x,6) {\n\\[0.3mm]\t};
    \node[align=center,font=\fontsize{6.4}{8}\selectfont,text=black!55] at (\x,-7) {\p};
  }
\end{tikzpicture}
\end{center}
\vspace{0.5mm}
\begin{columns}[T,onlytextwidth]
\begin{column}{0.50\textwidth}
  \begin{keybox}{全课主线}
  \small
  燃料化学能 $\rightarrow$ 锅炉产生蒸汽 $\rightarrow$ 汽轮机（或燃气轮机、内燃机）输出机械功 $\rightarrow$ 发电机输出电能。
  \par\vspace{1mm}
  \textbf{三个反复出现的问题}：能量边界在哪里？介质往哪里流？公式在什么假设下成立？
  \end{keybox}
\end{column}
\begin{column}{0.47\textwidth}
  \begin{bluebox}{本书的“读图”方法}
  \small
  \begin{enumerate}
    \item 先找\textbf{闭合回路}或\textbf{主气流通道}。
    \item 再标\textbf{压力升/降}和\textbf{温度升/降}。
    \item 最后把设备、状态点和公式一一对应。
  \end{enumerate}
  \end{bluebox}
\end{column}
\end{columns}
\vspace{1mm}
\begin{plainbox}
\small
\textbf{符号总则}：行内公式用 \texttt{\$...\$}；行间公式单独占行；\(\mathrm{d}\) 表示微分；单位用 \(\mathrm{MPa}\)、\(\mathrm{kW}\)、\(\mathrm{kJ/kg}\) 等正体；中文解释放在公式外。
\end{plainbox}
\end{frame}
